\chapter*{Abstract}
\addcontentsline{toc}{chapter}{Abstract}

High-Performance Computing (HPC) clusters have become essential infrastructure for scientific research, engineering simulations, and large-scale data processing. The Slurm Workload Manager is widely used across HPC environments and provides sophisticated resource management capabilities, but its command-line interface can create usability barriers for users who need to inspect queues, diagnose jobs, submit workloads, or perform controlled state-changing operations. This thesis addresses these challenges by designing and implementing an intelligent agent system that provides a natural language interface for Slurm cluster management.

The proposed Slurm Agent is a domain-specific control and assistance system for translating natural-language intent into concrete Slurm operations. Rather than treating an LLM as a general chatbot, the system combines typed Slurm tools, an Observer/Operator workflow with human confirmation for state-changing actions, local Slurm documentation retrieval, stateful confirmation handling, and trace-oriented evaluation. The implementation uses a React + Vite presentation layer, a session-aware FastAPI agent backend, an MCP protocol/tool layer, and real or mock Slurm infrastructure, but these frameworks serve as the substrate for Slurm-specific safety, grounding, and evaluation mechanisms.

This thesis makes three main contributions. \textbf{(1)} It introduces a capacity-aware dual-agent architecture that structurally separates read-only observation from state-changing execution: the Observer agent accesses 29 read-only tools while the Operator agent holds 40 action tools behind a mandatory human-in-the-loop confirmation gate, ensuring that destructive operations cannot execute without explicit user approval regardless of model behaviour. Ablation experiments confirm a +6.1\,pp overall improvement and +10.2\,pp on routing-neutral cases compared to a monolithic 64-tool baseline, attributable to tool-scope reduction alone. \textbf{(2)} It contributes a 3,135-case curated Slurm benchmark with architecture-aware scoring: each case carries ground-truth labels for expected tools, routing decisions, HITL requirements, and source-to-target state transitions, enabling evaluation at the behavioral level rather than text-quality level. The benchmark covers 11 task categories across 5 mock cluster scenarios and is publicly released as an independently reusable artifact. \textbf{(3)} It develops a role-scoped fine-tuning methodology that distills a commercial model's behaviour into an open-weight model (Qwen2.5-14B-Instruct via QLoRA): training traces are split at handoff boundaries and each sample's declared tool set is restricted to the active agent role. The resulting model achieves $88.3\% \pm 2.6$\,pp (95\% CI) mean weighted score, closing 29\% of the gap to the commercial baseline while running entirely on local infrastructure.

The implementation demonstrates the feasibility of a stateful, tool-grounded natural language interface for HPC operations. The fine-tuned open-weight model achieves an 88.3\% mean weighted score on the held-out test split, closing 29\% of the gap to the commercial API baseline while eliminating external API dependency. Domain-adapted open-weight models serve as viable local alternatives for structured tool-calling tasks, reducing operational cost and preserving data sovereignty. The system provides a practical foundation for AI-assisted cluster operations where usability, traceability, and controlled execution must be balanced.

\textbf{Keywords:} High-Performance Computing, Slurm, Large Language Models, Model Context Protocol, Multi-Agent Systems, Human-in-the-Loop, Natural Language Interface, Parameter-Efficient Fine-Tuning
